\section{How to use this file}
\label{sec:org2872d7b}
Plain lists are the lists you write in the body of an entry: bullets,
numbers, descriptions and checkboxes. Org understands their structure, so
it can renumber them, move items with their sub-items, and count the
checked boxes for you. This file goes from simple lists to checkbox
statistics.

Start Neovim from the root of the repository with the bundled config:

\begin{verbatim}
nvim -u examples/minimal_init.lua examples/03-lists.org
\end{verbatim}

\begin{itemize}
\item The file opens folded. \texttt{<Tab>} on a heading opens it, \texttt{<S-Tab>} cycles
the whole buffer.
\item \texttt{<prefix>} means \texttt{<leader>o} (\texttt{<Space>o} with the bundled config). \texttt{g?}
lists every key of the buffer.
\item \texttt{u} undoes; \texttt{git checkout examples/03-lists.org} restores the file.
\item \textbf{Try:} lines are exercises, \textbf{Expect:} lines say what you should see.
\item Lines starting with =\# = are comments: notes for you. A comment line in
column 0 ends a list, so comments are placed above or below lists, never
between their items.
\item The bundled config draws checkboxes as icons: \texttt{[ ]} as an empty box,
\texttt{[-]} as \texttt{◐} and \texttt{[X]} as \texttt{✓}. The text in the file does not change.
\item "On an item" means the cursor is on the first line of the item, usually
on its text (not in column 0).
\end{itemize}
\subsection{Keys in this file}
\label{sec:orgc60e13e}
\begin{center}
\begin{tabular}{ll}
Key & What it does\\
\hline
\texttt{<M-CR>} & new item (Normal mode: after this one)\\
\texttt{<M-S-CR>} & new item with a checkbox\\
\texttt{<prefix>is} & new sub-item (indented)\\
\texttt{<Tab>} (Insert mode) & on an empty item: cycle its indentation\\
\texttt{<Tab>} & on an item with sub-items: fold / unfold\\
\texttt{<S-Right>} \texttt{<S-Left>} & cycle the bullet style of the list\\
\texttt{<prefix>hb} / \texttt{<C-c>-} & cycle the bullet style\\
\texttt{<S-Down>} \texttt{<S-Up>} & next / previous item of the same level\\
\texttt{>{}>{}} \texttt{<{}<{}} / \texttt{<M-l>} \texttt{<M-h>} & indent / outdent the item\\
\texttt{>s} \texttt{<s} / \texttt{<M-L>} \texttt{<M-H>} & indent / outdent the item and its children\\
\texttt{<M-k>} \texttt{<M-j>} & move the item (with sub-items) up / down\\
\texttt{<C-Space>} & toggle a checkbox\\
\texttt{<C-c><C-c>} & toggle a checkbox, or repair the list\\
\texttt{<C-c><C-x><C-b>} & toggle checkbox (Emacs key)\\
\texttt{<prefix>\#} / \texttt{<C-c>\#} & update statistics cookies\\
\texttt{<prefix>-} & toggle between item and text / headline\\
\texttt{<prefix>*} & turn the item into a headline\\
\texttt{<C-c><C-*>} & turn the whole list into a subtree\\
\texttt{<prefix>hs} & sort the list\\
\texttt{<C-c><C-x><C-r>} & toggle a radio button\\
\end{tabular}
\end{center}
\section{Unordered lists}
\label{sec:org27a19fc}
An item starts with a bullet followed by a space: \texttt{-}, \texttt{+}, or \texttt{*} (a star
only when the item is indented, because a star in column 0 starts a
headline). Lines that belong to an item are indented past its bullet.
Items indented under another item form a sub-list.

A list ends:
\begin{itemize}
\item at a line that is indented at or left of the first bullet (and is not a
new item),
\item at a headline,
\item or after \textbf{two} blank lines in a row.
\end{itemize}

\textbf{Try:} open "Groceries" below. Put the cursor on "Fruit" and press \texttt{<Tab>},
then \texttt{<Tab>} again.

\textbf{Expect:} the first \texttt{<Tab>} folds "Fruit": its sub-items "apples" and
"pears" and its second line are hidden, the line ends with \texttt{...}. The
second \texttt{<Tab>} shows them again. (\texttt{cycle\_include\_plain\_lists = true}, the
default, makes items with sub-items or several lines foldable.)
\subsection{Groceries}
\label{sec:orgaf6c356}
\begin{itemize}
\item Fruit
(only the ripe ones)
\begin{itemize}
\item apples
\item pears
\end{itemize}
\item Vegetables
\begin{itemize}
\item carrots
\item leeks
\begin{itemize}
\item the thin ones
\end{itemize}
\end{itemize}
\item Bread
\end{itemize}
This line is not indented, so it is not part of the list.
\subsection{Where a list ends}
\label{sec:orge1f6803}
\begin{itemize}
\item first item
\item second item
\end{itemize}


This paragraph comes after two blank lines: the list above has ended.

\begin{verbatim}
- this looks like an item
- but it is example text
\end{verbatim}
\section{Ordered lists}
\label{sec:org9627784}
Ordered bullets are a number followed by \texttt{.} or \texttt{)}: \texttt{1.}, \texttt{1)}. You never
renumber by hand: every list command renumbers the list when it changes it.
If the numbers are wrong (you typed them, or pasted lines), \texttt{<C-c><C-c>} on
the first line of an item repairs the whole list.

A counter \texttt{[@N]} right after the bullet makes the list continue from N.

\textbf{Try:} put the cursor on "Wake up" (in "Wrong numbers") and press
\texttt{<C-c><C-c>}.

\textbf{Expect:} the numbers become 1, 2, 3, 4:
\begin{verbatim}
1. Wake up
2. Make coffee
3. Drink coffee
4. Work
\end{verbatim}

\textbf{Try:} in "Counter", press \texttt{<C-c><C-c>} on "Chapter five".

\textbf{Expect:} the numbers stay 5, 6, 7: the list starts at the counter.
Change \texttt{[@5]} to \texttt{[@10]} and press \texttt{<C-c><C-c>} again: 10, 11, 12.
\subsection{Wrong numbers}
\label{sec:org3030266}
\begin{enumerate}
\item Wake up
\item Make coffee
\item Drink coffee
\item Work
\end{enumerate}
\subsection{Counter}
\label{sec:org10c0f82}
\begin{enumerate}
\setcounter{enumi}{4}
\item Chapter five
\item Chapter six
\item Chapter seven
\end{enumerate}
\subsection{Two styles}
\label{sec:org8463113}
\begin{enumerate}
\item A closing parenthesis
\item works the same way
\begin{enumerate}
\item and a sub-list can use its own style
\item independently
\end{enumerate}
\end{enumerate}
\section{Description lists}
\label{sec:orge8718bb}
A description item has a term, then \texttt{::} surrounded by spaces, then the
description. It is an unordered item (\texttt{-} or \texttt{+}), and it is exported as a
definition list (\texttt{<dl>} in HTML).

\textbf{Try:} on "Neovim" press \texttt{\$} (end of line) then \texttt{<M-CR>}, type \texttt{Emacs},
press \texttt{<Esc>}.

\textbf{Expect:} a new description item right after "Neovim", with the separator
already there and the cursor on the term:
\begin{verbatim}
- Emacs ::
\end{verbatim}
\subsection{Editors}
\label{sec:orgc5d4884}
\begin{description}
\item[{Org}] a plain-text outliner and organizer
\item[{Neovim}] a hyperextensible Vim-based text editor
\item[{Plain text}] lasts forever
and a description can go on over
several lines.
\end{description}
\section{Inserting items}
\label{sec:orgbdb4163}
\texttt{<M-CR>} (Meta-Return) on an item inserts a new item of the same kind
(bullet, description term) and leaves you in Insert mode. Where:
\begin{itemize}
\item Normal mode, anywhere on the item's first line (even column 0): after
the item \textbf{and its sub-items}.
\item Insert mode: at the cursor. At or before the item text (e.g. \texttt{i} in
column 0) the new item goes above; in the middle of the text, the rest of
the text moves to the new item.
\end{itemize}

\texttt{<M-S-CR>} does the same but adds an empty checkbox \texttt{[ ]}. \texttt{<prefix>is}
inserts a sub-item (indented one level). A new item never copies the
checkbox of the current one with \texttt{<M-CR>}.

In Insert mode, right after \texttt{<M-CR>}, \texttt{<Tab>} on the new empty item cycles
its indentation: under the previous item, then back out level by level.

When the items of a list are separated by blank lines, new items get a
blank line too (\texttt{blank\_before\_new\_entry = \{ plain\_list\_item = "auto" \}}).

\textbf{Try:} on "Saturday" press \texttt{\$} and \texttt{<M-CR>}, type \texttt{Sunday}, \texttt{<Esc>}.

\textbf{Expect:} \texttt{- Sunday} after "Saturday" (and after its sub-item "morning
run"), before "Next week".

\textbf{Try:} on "Saturday" press \texttt{\$} and \texttt{<M-S-CR>}, type \texttt{Pack}, \texttt{<Esc>}.

\textbf{Expect:} \texttt{- [ ] Pack} after the "Saturday" item.

\textbf{Try:} on "Saturday" press \texttt{\$} and \texttt{<prefix>is}, type \texttt{evening}, \texttt{<Esc>}.

\textbf{Expect:} an indented sub-item =  - evening= right below "Saturday".

\textbf{Try:} on "Next week" press \texttt{A}, then \texttt{<M-CR>}, then \texttt{<Tab>}, type
\texttt{Monday} and \texttt{<Esc>}.

\textbf{Expect:} =  - Monday= indented under "Next week".

\textbf{Try:} on "Split this item" put the cursor on \texttt{i} of "item", press \texttt{i} and
\texttt{<M-CR>}, then \texttt{<Esc>}.

\textbf{Expect:} two items: \texttt{- Split this} and \texttt{- item}.
\subsection{Weekend}
\label{sec:orge87d85b}
\begin{itemize}
\item Saturday
\begin{itemize}
\item morning run
\end{itemize}
\item Next week
\item Split this item
\end{itemize}
\subsection{Spaced items}
\label{sec:org99371d7}
\begin{itemize}
\item first

\item second
\end{itemize}

\textbf{Try:} on "second" press \texttt{\$} and \texttt{<M-CR>}, type \texttt{third}.

\textbf{Expect:} a blank line, then \texttt{- third}.
\section{Bullet styles}
\label{sec:org1183b91}
\texttt{<S-Right>} / \texttt{<S-Left>} on an item cycle the bullet of \textbf{the whole list}
(the item and its siblings) through \texttt{-}, \texttt{+}, \texttt{*}, \texttt{1.}, \texttt{1)}. The star is
skipped for a list in column 0. Sub-lists keep their own bullets.
\texttt{<prefix>hb} and \texttt{<C-c>-} cycle forward too.

\textbf{Try:} on "alpha" (in "Cycle me") press \texttt{<S-Right>} four times.

\textbf{Expect:} the bullets of alpha, beta and gamma become \texttt{+}, then \texttt{1.} \texttt{2.}
\texttt{3.}, then \texttt{1)} \texttt{2)} \texttt{3)}, then \texttt{-} again. The sub-item "beta one" keeps its
\texttt{-} bullet (and moves right when the numbers make the bullet wider).
\subsection{Cycle me}
\label{sec:org6852f3b}
\begin{itemize}
\item alpha
\item beta
\begin{itemize}
\item beta one
\end{itemize}
\item gamma
\end{itemize}
\section{Moving between items}
\label{sec:orgc4eb537}
\texttt{<S-Down>} / \texttt{<S-Up>} jump to the next / previous item \textbf{of the same level},
skipping sub-items.

\textbf{Try:} on "Ann" press \texttt{<S-Down>} twice.

\textbf{Expect:} the cursor goes to "Bob" (skipping "Ann's notes"), then to "Cid".
\subsection{People}
\label{sec:org1a9b2bc}
\begin{itemize}
\item Ann
\begin{itemize}
\item Ann's notes
\end{itemize}
\item Bob
\item Cid
\end{itemize}
\section{Indenting and outdenting}
\label{sec:org42f874c}
\begin{center}
\begin{tabular}{ll}
Key & Moves\\
\hline
\texttt{>{}>{}} / \texttt{<{}<{}}, \texttt{<M-l>} / \texttt{<M-h>} & the item alone\\
\texttt{>s} / \texttt{<s}, \texttt{<M-L>} / \texttt{<M-H>} & the item and its sub-items (subtree)\\
\end{tabular}
\end{center}

Org's list rules apply:
\begin{itemize}
\item The first item of a list cannot be indented alone; \texttt{<M-L>} / \texttt{<M-H>} on
it move the whole list.
\item An item with sub-items cannot be outdented alone: use \texttt{<s} or \texttt{<M-H>}.
\item After each change the list is repaired: numbering, bullets, the
indentation of sub-lists and parent checkboxes.
\item In Visual mode, \texttt{<M-l>} / \texttt{<M-h>} indent / outdent every selected item.
\end{itemize}

\textbf{Try:} on "Task B" press \texttt{>s}.

\textbf{Expect:} "Task B" and its sub-item "B detail" move right together; "Task
B" is now a sub-item of "Task A":
\begin{verbatim}
- Task A
  - Task B
    - B detail
- Task C
\end{verbatim}

\textbf{Try:} press \texttt{u}, then \texttt{>{}>{}} on "Task B".

\textbf{Expect:} only "Task B" moves; "B detail" stays where it was and becomes
its sibling:
\begin{verbatim}
- Task A
  - Task B
  - B detail
\end{verbatim}

\textbf{Try:} press \texttt{u}, then \texttt{>{}>{}} on "Task A".

\textbf{Expect:} nothing changes and a message says "At first item: use
S-M-<left/right> to move the whole list".

\textbf{Try:} on "Nested parent" (in "Outdent rules") press \texttt{<{}<{}}.

\textbf{Expect:} the message "Cannot outdent an item without its children".
Press \texttt{<s} instead: "Nested parent" and "Nested child" both move left.
\subsection{Indent rules}
\label{sec:orgffe733b}
\begin{itemize}
\item Task A
\item Task B
\begin{itemize}
\item B detail
\end{itemize}
\item Task C
\end{itemize}
\subsection{Outdent rules}
\label{sec:org765f93b}
\begin{itemize}
\item Top
\begin{itemize}
\item Nested parent
\begin{itemize}
\item Nested child
\end{itemize}
\end{itemize}
\end{itemize}
\section{Moving items}
\label{sec:org4f04a7e}
\texttt{<M-k>} / \texttt{<M-j>} (also \texttt{<M-Up>} / \texttt{<M-Down>}) move an item up / down past
its sibling. Sub-items, extra lines and checkboxes travel with it, and
ordered lists are renumbered.

\textbf{Try:} on "Third" press \texttt{<M-k>} twice.

\textbf{Expect:}
\begin{verbatim}
1. Third
2. First
3. Second
\end{verbatim}
The numbers stay in order; only the texts moved.

\textbf{Try:} on "Unpacked" press \texttt{<M-j>}.

\textbf{Expect:} "Unpacked" moves below "Packed" together with its sub-item:
\begin{verbatim}
- [X] Packed
- [ ] Unpacked
  - [ ] socks
\end{verbatim}
\subsection{Numbered to reorder}
\label{sec:org4f0e106}
\begin{enumerate}
\item First
\item Second
\item Third
\end{enumerate}
\subsection{Boxes to reorder}
\label{sec:orgd2d6e17}
\begin{itemize}
\item[{$\square$}] Unpacked
\begin{itemize}
\item[{$\square$}] socks
\end{itemize}
\item[{$\boxtimes$}] Packed
\end{itemize}
\section{Converting items, text and headlines}
\label{sec:org15303e2}
\begin{center}
\begin{tabular}{lll}
Key & On & Result\\
\hline
\texttt{<prefix>-} & a text line & an item\\
\texttt{<prefix>-} & an item & a text line (the checkbox stays)\\
\texttt{<prefix>-} & a headline & an item (a TODO keyword becomes a box)\\
\texttt{<prefix>*} & an item & a headline (\texttt{[ ]} / \texttt{[X]} → TODO / DONE)\\
\texttt{<C-c><C-*>} & an item & the whole list becomes a subtree\\
\end{tabular}
\end{center}

In Visual mode \texttt{<prefix>-} converts every selected line; with a count the
selection becomes a single item.

\textbf{Try:} select the three lines "eggs", "flour" and "sugar" (in "Text to
items") with \texttt{V} and \texttt{2j}, then press \texttt{<prefix>-}.

\textbf{Expect:}
\begin{verbatim}
- eggs
- flour
- sugar
\end{verbatim}

\textbf{Try:} in "Items to headlines", on "[ ] Book the venue" press \texttt{<C-c><C-*>}.

\textbf{Expect:} the list is gone; its items are now child headlines of "Items to
headlines", with checkboxes turned into keywords and sub-items into deeper
headlines:
\begin{verbatim}
*** TODO Book the venue
*** DONE Send the invitations
**** DONE Family
**** DONE Friends
\end{verbatim}
\subsection{Text to items}
\label{sec:org38bae4b}
eggs
flour
sugar
\subsection{Items to headlines}
\label{sec:orgefbd59a}
\begin{itemize}
\item[{$\square$}] Book the venue
\item[{$\boxtimes$}] Send the invitations
\begin{itemize}
\item[{$\boxtimes$}] Family
\item[{$\boxtimes$}] Friends
\end{itemize}
\end{itemize}
\section{Checkboxes}
\label{sec:org80a4529}
An item becomes a task when it starts with a checkbox: \texttt{[ ]} (open),
\texttt{[X]} (done) or \texttt{[-]} (partly done, set automatically).

\begin{itemize}
\item \texttt{<C-Space>} (or \texttt{<C-c><C-c>}, or Emacs \texttt{<C-c><C-x><C-b>}) on the first
line of an item toggles its box.
\item A parent box follows its children: when some children are done it shows
\texttt{[-]}, when all are done \texttt{[X]}. Toggling the parent itself is refused
("Cannot toggle this checkbox").
\item With a count: \texttt{4<C-c><C-c>} adds (or removes) a checkbox on the item;
\texttt{16<C-c><C-c>} sets \texttt{[-]}.
\item In Visual mode, \texttt{<C-Space>} toggles every selected item (following the
state of the first one).
\item On a \textbf{headline}, \texttt{<C-Space>} toggles every item of its text.
\end{itemize}

\textbf{Try:} in "Packing", toggle "Passport" and "Tickets" with \texttt{<C-Space>}.

\textbf{Expect:} "Documents" first becomes \texttt{[-]} (after Passport), then \texttt{[X]}
(after Tickets). "Documents" was never toggled by you.

\textbf{Try:} on "Documents" press \texttt{<C-c><C-c>}.

\textbf{Expect:} nothing changes, and the message "Cannot toggle this checkbox:
all subitems checked" appears. (\texttt{<C-Space>} silently leaves it alone.)

\textbf{Try:} on "Toothbrush" (no box yet) press \texttt{4<C-c><C-c>}.

\textbf{Expect:} \texttt{- [ ] Toothbrush}. Press \texttt{4<C-c><C-c>} again to remove it.

\textbf{Try:} select "Shirts", "Socks" and "Shoes" with \texttt{V2j} and press
\texttt{<C-Space>}.

\textbf{Expect:} all three become \texttt{[X]}.
\subsection{Packing}
\label{sec:orgbdd1cd6}
\begin{itemize}
\item[{$\square$}] Documents
\begin{itemize}
\item[{$\square$}] Passport
\item[{$\square$}] Tickets
\end{itemize}
\item Toothbrush
\item[{$\square$}] Shirts
\item[{$\square$}] Socks
\item[{$\square$}] Shoes
\end{itemize}
\subsection{Toggle everything from the headline}
\label{sec:org3ec8003}
\begin{itemize}
\item[{$\square$}] one
\item[{$\square$}] two
\item[{$\square$}] three
\end{itemize}

\textbf{Try:} on the headline "Toggle everything from the headline" press
\texttt{<C-Space>}.

\textbf{Expect:} one, two and three are all checked. Press it again to uncheck
them.
\subsection{Ordered checklists}
\label{sec:org302e170}
With the property \texttt{ORDERED: t}, the boxes must be checked in order: an
unchecked box blocks the ones after it.

\textbf{Try:} press \texttt{<C-c><C-c>} on "Step two" (before "Step one").

\textbf{Expect:} "Step two" stays unchecked and the message "Cannot toggle this
checkbox: unchecked subitems" appears (the same words as Emacs). Check
"Step one" first; then "Step two" can be checked.

\begin{itemize}
\item[{$\square$}] Step one
\item[{$\square$}] Step two
\item[{$\square$}] Step three
\end{itemize}
\section{Statistics cookies}
\label{sec:org45cb940}
A cookie \texttt{[/]} or \texttt{[\%]} at the end of an item or a headline shows how many
of its checkboxes are done: \texttt{[2/5]} or \texttt{[40\%]}. Type the empty cookie
yourself; Org fills it in and updates it after every toggle. You can put
both on the same line.

\begin{itemize}
\item On an \textbf{item}, the cookie counts its direct sub-items.
\item On a \textbf{headline}, it counts the top-level items of its text (sub-items are
counted through their parents), or, when there are no checkboxes, its
child TODO entries (see \href{04-todo.org}{04-todo.org}).
\item \texttt{<prefix>\#} (Emacs \texttt{<C-c>\#}) updates the cookies of the current entry;
\texttt{4<prefix>\#} updates every cookie of the buffer.
\end{itemize}

\textbf{Try:} go to "Party [/] [\%]" below and press \texttt{<prefix>\#} on its headline.

\textbf{Expect:} the headline shows \texttt{[1/3] [33\%]} (one of its three top-level
items, "Music", is checked). "Food" shows \texttt{[1/2]}.

\textbf{Try:} check "Snacks".

\textbf{Expect:} "Food" becomes \texttt{[X]} with \texttt{[2/2]}, and the headline becomes
\texttt{[2/3] [66\%]}.
\subsection{Party [/] [\%]}
\label{sec:org8242909}
\begin{itemize}
\item[{$\boxminus$}] Food [/]
\begin{itemize}
\item[{$\boxtimes$}] Cake
\item[{$\square$}] Snacks
\end{itemize}
\item[{$\boxtimes$}] Music
\item[{$\square$}] Invitations
\end{itemize}
\subsection{Counting the whole tree [/]}
\label{sec:orgde2857d}
\texttt{COOKIE\_DATA} changes what a headline cookie counts: \texttt{checkbox} or \texttt{todo}
forces one kind, and \texttt{recursive} counts every box in the tree, sub-items
included.

\textbf{Try:} press \texttt{<prefix>\#} on this headline.

\textbf{Expect:} \texttt{[2/7]}: all seven boxes are counted, parents and sub-items
alike (Food, Cake, Snacks, Decoration, Balloons, Lights, Music), and two of
them (Cake and Music) are checked. Without \texttt{recursive} the cookie would
count only the three top-level items: \texttt{[1/3]}.

\begin{itemize}
\item[{$\boxminus$}] Food
\begin{itemize}
\item[{$\boxtimes$}] Cake
\item[{$\square$}] Snacks
\end{itemize}
\item[{$\square$}] Decoration
\begin{itemize}
\item[{$\square$}] Balloons
\item[{$\square$}] Lights
\end{itemize}
\item[{$\boxtimes$}] Music
\end{itemize}
\subsection{Cookies without checkboxes}
\label{sec:orge6bfee2}
When a headline has a cookie but no checkboxes in its text and no child
entries, \texttt{<prefix>\#} sets it to \texttt{[0/0]} or \texttt{[100\%]}.
\subsubsection{Nothing to count [\%]}
\label{sec:orgb08c2e4}
\textbf{Try:} press \texttt{<prefix>\#} on "Nothing to count".

\textbf{Expect:} \texttt{[100\%]}.
\section{Radio lists}
\label{sec:orgb72b707}
A radio list allows only one checked item: checking one unchecks the
others. Mark the list with \texttt{\#+attr\_org: :radio t} on the line above it;
then \texttt{<C-c><C-c>} and \texttt{<C-Space>} work like radio buttons.
\texttt{<C-c><C-x><C-r>} (\texttt{:Org toggle\_radio\_button}) toggles a radio button in
any list, and \texttt{:Org checkbox\_radio\_mode} makes \texttt{<C-c><C-c>} behave that
way in every list of the buffer.

\textbf{Try:} in "Coffee size", press \texttt{<C-c><C-c>} on "Large".

\textbf{Expect:} "Large" is checked and "Medium" is unchecked.

\textbf{Try:} in "Normal list", put the cursor on "red" and press
\texttt{<C-c><C-x><C-r>}.

\textbf{Expect:} "red" is checked and "green" unchecked, although this list has no
\texttt{:radio} attribute.
\subsection{Coffee size}
\label{sec:org0c92d11}
\begin{itemize}
\item[{$\square$}] Small
\item[{$\boxtimes$}] Medium
\item[{$\square$}] Large
\end{itemize}
\subsection{Normal list}
\label{sec:orgd57504c}
\begin{itemize}
\item[{$\square$}] red
\item[{$\boxtimes$}] green
\item[{$\square$}] blue
\end{itemize}
\section{Sorting lists}
\label{sec:orgb0f1a4e}
\texttt{<prefix>hs} (Emacs \texttt{<C-c>\textasciicircum{}}) on an item sorts that item and its siblings.
The menu has fewer keys than for headlines:

\begin{center}
\begin{tabular}{ll}
Key & Sorts by\\
\hline
\texttt{a} & the item text, alphabetically\\
\texttt{n} & the number at the start of the text\\
\texttt{t} & the first timestamp of the item (or a timer \texttt{0:05:00 ::} term)\\
\texttt{x} & checkbox: unchecked \texttt{[ ]} first, then \texttt{[-]}, then \texttt{[X]}\\
\texttt{f} & a Lua function\\
\end{tabular}
\end{center}

Uppercase keys reverse the order; a count makes \texttt{a} case-sensitive.
Sub-items travel with their item, and ordered lists are renumbered.

\textbf{Try:} on "cherry" press \texttt{<prefix>hs} then \texttt{a}.

\textbf{Expect:} apple, banana (with its sub-item), cherry, numbered 1, 2, 3.

\textbf{Try:} on any item of "Chores list" press \texttt{<prefix>hs} then \texttt{x}.

\textbf{Expect:} the open items first, each group in its original order:
\begin{verbatim}
- [ ] vacuum
- [ ] windows
- [X] dishes
- [X] laundry
\end{verbatim}
Press \texttt{<prefix>hs} then \texttt{X} (reverse) to get the checked ones first.

\textbf{Try:} on "Retro" press \texttt{<prefix>hs} then \texttt{t}.

\textbf{Expect:} Kickoff (Oct 1), Review (Oct 14), Retro (Nov 2).
\subsection{Fruit list}
\label{sec:orgf7915eb}
\begin{enumerate}
\item cherry
\item banana
\begin{itemize}
\item a sub-item stays with banana
\end{itemize}
\item apple
\end{enumerate}
\subsection{Chores list}
\label{sec:org99af24e}
\begin{itemize}
\item[{$\square$}] vacuum
\item[{$\boxtimes$}] dishes
\item[{$\square$}] windows
\item[{$\boxtimes$}] laundry
\end{itemize}
\subsection{Meetings list}
\label{sec:org9b58fed}
\begin{itemize}
\item Retro \textit{<2026-11-02 Mon>}
\item Meeting: Kickoff \textit{<2026-10-01 Thu>}
\item Review \textit{<2026-10-14 Wed>}
\end{itemize}
\section{Checkboxes that block a TODO}
\label{sec:orgd39765f}
With \texttt{enforce\_todo\_checkbox\_dependencies = true}, an entry cannot be marked
DONE while it has unchecked boxes. It is off by default; for this session:

\begin{verbatim}
:lua require("org.config").opts.enforce_todo_checkbox_dependencies = true
\end{verbatim}

\textbf{Try:} turn it on, then on "Ship the release" press \texttt{cit} (next TODO
state).

\textbf{Expect:} the entry stays TODO and a message says it is blocked. Check both
boxes and press \texttt{cit} again: it becomes DONE.
\subsection{{\bfseries\sffamily TODO} Ship the release}
\label{sec:orgfe04580}
\begin{itemize}
\item[{$\square$}] Tests pass
\item[{$\square$}] Changelog written
\end{itemize}
\section{List settings}
\label{sec:org6d2fcfd}
\begin{center}
\begin{tabular}{ll}
Option (default) & Effect\\
\hline
\texttt{blank\_before\_new\_entry.plain\_list\_item} (auto) & blank lines between items\\
\texttt{cycle\_include\_plain\_lists} (\texttt{true}) & \texttt{<Tab>} folds items\\
\texttt{meta\_return\_split\_line} (\texttt{true}) & \texttt{<M-CR>} splits the text\\
\texttt{enforce\_todo\_checkbox\_dependencies} (\texttt{false}) & open boxes block DONE\\
\texttt{ui.checkboxes} (\texttt{false}) & icons for the boxes\\
\end{tabular}
\end{center}

Not available (Emacs options without an org.nvim equivalent): alphabetical
bullets (\texttt{a.}, \texttt{org-list-allow-alphabetical}), changing the bullet
automatically when demoting (\texttt{org-list-demote-modify-bullet}) and a custom
indentation of sub-lists (\texttt{org-list-indent-offset}).

Timer lists (items whose term is a running time, \texttt{0:05:12 ::}) are covered
in \href{20-timers-reminders.org}{20-timers-reminders.org}.
\section{Further reading}
\label{sec:orgd5b1c41}
\begin{description}
\item[{\texttt{:h org-lists}}] lists, checkboxes and cookies
\item[{\texttt{:h org-radio-list}}] radio buttons
\item[{\texttt{:h org-meta-return}}] \texttt{<M-CR>} in lists
\item[{\texttt{:h org-promote-demote}}] indentation rules for items
\item[{\texttt{:h org-sort}}] sorting lists
\item[{\texttt{:h org-todo-statistics}}] cookies that count TODO entries
\end{description}
